\subsection{Non-Committing Encryption}
\label{sec:nce}

Non-committing encryption (NCE) was introduced by Canetti et al.\
\cite{CanettiEtAl1996AdaptiveMPC}.  It lets a simulator prepare a public
key and ciphertext without knowing the message, and later explain the
sender's and receiver's random tapes consistently with that message.
We give the key-first definition of Hemenway et al.\
\cite[Definition~2]{HemenwayEtAl2016QuasiOptimalNCE} explicit game
commands; the same definition appears in Brakerski et al.\
\cite[full version, Definition~11]{BrakerskiEtAl2020ConstantNCE}.
In particular, the message may be chosen after the public key is known.
Full random-tape exposure, rather than exposure of the final secret key
alone, is essential when there are no erasures; see also
Damg{\aa}rd and Nielsen~
\cite[full version, Section~4]{DamgardNielsen2000NCE}.

\subsubsection{Syntax and correctness}
Fix a public message length $\ell=\ell(\lambda)$, polynomial in the
security parameter.  All algorithms below are parameterized by $\ell$;
we omit this fixed parameter from their calls.  The encryption scheme
has polynomial-time algorithms
\[
  (\pk,\sk)\sample\KeyGen(1^\lambda),\qquad
  c\sample\Enc(\pk,m),\qquad
  m'\sample\Dec(\sk,c),
\]
where $m\in\bits^\ell$ and decryption returns an element of
$\bits^\ell\cup\{\bot\}$.  We allow randomized decryption, as in the
constant-rate construction cited below.  When its coins matter we write
$\Dec(\sk,c;\delta)$; otherwise fresh independent coins are understood.
Keys and ciphertexts have prescribed polynomial lengths, using padding
if needed.
Correctness requires, for every efficient algorithm $\mathcal B$ that
chooses a message $m\in\bits^\ell$ after seeing $\pk$,
\[
  \Pr\bigl[\Dec(\sk,c)=m\bigr]\geq 1-\negl(\lambda),
\]
where $(\pk,\sk)\sample\KeyGen(1^\lambda)$,
$m\sample\mathcal B(1^\lambda,\pk)$, and
$c\sample\Enc(\pk,m)$.  The probability includes the decryption coins.

Write $\rho_R$ and $\rho_S$ for the \emph{complete} uniform random
tapes of key generation and encryption, respectively, each of a fixed
polynomial length.  Thus
\[
  (\pk,\sk)=\KeyGen(1^\lambda;\rho_R),\qquad
  c=\Enc(\pk,m;\rho_S)
\]
determine the two computations completely.  A receiver that subsequently
decrypts also retains its decryption coins; these are fresh local coins,
not part of $\rho_R$.  The experiment below has no private decryption
step before corruption.

\subsubsection{Simulation with random-tape exposure}
A simulator consists of two efficient algorithms:
\begin{description}
  \item[$\SimNCE(1^\lambda)\to(\widetilde{\pk},\widetilde c,\SimState)$.]
    Prepare a simulated public key and ciphertext, retaining private
    state $\SimState$, without receiving the message.
  \item[$\ExplainNCE(\SimState,m)\to
      (\widetilde\rho_R,\widetilde\rho_S)$.]
    Given the message, explain the complete key-generation and
    encryption tapes for the previously prepared pair.
\end{description}

The experiments $\NCEReal_{\mathcal A}(1^\lambda)$ and
$\NCEIdeal_{\mathcal A,\mathcal S}(1^\lambda)$ give $1^\lambda$ and
$\ell$ to the efficient adversary $\mathcal A$ and offer these commands:
\begin{description}
  \item[$\GetPublicKey()$.]
    In the real experiment, sample $\rho_R$ uniformly and compute
    \[
      (\pk,\sk)\gets\KeyGen(1^\lambda;\rho_R).
    \]
    Retain $(\sk,\rho_R)$ and return $\pk$.
    In the ideal experiment, compute
    \[
      (\widetilde{\pk},\widetilde c,\SimState)\sample\SimNCE(1^\lambda),
    \]
    retain $(\widetilde c,\SimState)$, and return $\widetilde{\pk}$.
    The precomputed ciphertext remains private until $\EncryptMessage$.
  \item[$\EncryptMessage(m)$.]
    Store $m\in\bits^\ell$, which may depend on the public key.
    In the real experiment, sample $\rho_S$ uniformly and return
    $c\gets\Enc(\pk,m;\rho_S)$.
    In the ideal experiment, return the stored ciphertext $\widetilde c$.
    The simulator is not given $m$ at this step.
  \item[$\CorruptSender()$.]
    In the real experiment, return $(m,\rho_S)$.
    In the ideal experiment, return $(m,\widetilde\rho_S)$, using the
    cached explanation described below.
  \item[$\CorruptReceiver()$.]
    In the real experiment, return $(\sk,\rho_R)$.
    In the ideal experiment, compute
    $(\pk',\widetilde{\sk})\gets
      \KeyGen(1^\lambda;\widetilde\rho_R)$
    from the cached explanation, and return
    $(\widetilde{\sk},\widetilde\rho_R)$.
  \item[$\Finish(b)$.]
    Terminate and output $b\in\bits$.
\end{description}
On the \emph{first} corruption command in the ideal experiment, give
$m$ to the simulator, compute
\[
  (\widetilde\rho_R,\widetilde\rho_S)
    \sample\ExplainNCE(\SimState,m),
\]
and cache this pair for both corruption commands.

The commands $\GetPublicKey$ and $\EncryptMessage$ are accepted once
each, in that order.  Each corruption command is accepted at most once,
only after $\EncryptMessage$, and the two corruptions may occur in either
order.  The adversary may omit either or both.  Malformed, out-of-order,
or repeated commands return $\bot$ without changing state.  $\Finish$
may be called at any time; termination without an output bit means
$\Finish(0)$.  Allowing the adversary to request the two parts of one
explanation at different times is just post-processing of the joint
exposure in the cited definitions.

\begin{definition}[Non-committing security]
\label{def:nce}
The scheme has non-committing security with tape exposure if there exists
an efficient simulator $\mathcal S=(\SimNCE,\ExplainNCE)$ such that,
for every efficient adversary $\mathcal A$, there is a negligible
function $\varepsilon_{\mathcal A}$ satisfying
\[
  \left|
    \Pr[\NCEReal_{\mathcal A}(1^\lambda)=1]
    -\Pr[\NCEIdeal_{\mathcal A,\mathcal S}(1^\lambda)=1]
  \right|\leq\varepsilon_{\mathcal A}(\lambda).
\]
\end{definition}

\begin{definition}[Sender and receiver restrictions]
\label{def:sender-nce}
\label{def:receiver-nce}
\emph{Sender non-committing encryption} (SNCE) requires the same
indistinguishability only for adversaries that never call
$\CorruptReceiver$.  \emph{Receiver non-committing encryption} (RNCE)
requires it only for adversaries that never call $\CorruptSender$.
Each restriction permits its own simulator.  Full NCE implies both by
discarding the unused part of the explanation.  The receiver-only
notion originates with Jarecki and Lysyanskaya~
\cite{JareckiLysyanskaya2000AdaptiveThreshold}.
\end{definition}

In the full game the explained tapes must, except with negligible
probability, satisfy
\[
  \KeyGen(1^\lambda;\widetilde\rho_R)
    =(\widetilde{\pk},\widetilde{\sk}),\qquad
  \Enc(\widetilde{\pk},m;\widetilde\rho_S)=\widetilde c.
\]
They must also have the correct joint computational distribution, not
merely satisfy these equations.  Decrypting $\widetilde c$ with the
explained secret key and fresh coins recovers $m$ except with negligible
probability.  Explained tapes describe honest computations; they need
not be the tapes actually used internally by the simulator.

This definition covers only one ciphertext per key.  We use a
separate key for each input provider's entire prescribed-length input,
and separate delivery keys for the receiver-only applications.
The full NCE constructions of Brakerski et al.\
\cite[full version, Section~6, Theorems~23--24]{BrakerskiEtAl2020ConstantNCE}
give constant ciphertext expansion under their DDH or LWE assumptions.
For long messages,
their ciphertext length is $O(\ell)+\poly(\lambda)$; this rate excludes
public-key size and does not assert linear total computation.
